\chapter{Materials and Methods}

\section{Neighbourhood Structures}

A neighbourhood $\mathcal{N}(\pi)$ is the set of all permutations reachable
from $\pi$ by applying a single elementary move. Three neighbourhoods were
implemented.

\paragraph{Transpose.} Swaps two adjacent elements at positions $i$ and $i+1$:
\[
    \text{transpose}(\pi, i) =
    (\pi_1, \ldots, \pi_{i-1},\, \pi_{i+1},\, \pi_i,\, \pi_{i+2}, \ldots, \pi_n).
\]
The neighbourhood has $n-1$ neighbours. The change in objective value reduces
to a single matrix lookup:
\begin{equation}
    \Delta_T(\pi, i) = c_{\pi_{i+1}\pi_i} - c_{\pi_i\pi_{i+1}},
    \label{eq:delta_transpose}
\end{equation}
so each evaluation costs $O(1)$ and a full neighbourhood scan costs $O(n)$.

\paragraph{Exchange.} Swaps two arbitrary elements at positions $i < j$:
\[
    \text{exchange}(\pi, i, j) =
    (\pi_1, \ldots, \pi_j, \ldots, \pi_i, \ldots, \pi_n).
\]
The neighbourhood has $n(n-1)/2$ neighbours. The incremental delta is:
\begin{equation}
    \Delta_X(\pi, i, j) =
    (c_{\pi_j\pi_i} - c_{\pi_i\pi_j})
    + \sum_{k=i+1}^{j-1}
    \bigl[(c_{\pi_j\pi_k} - c_{\pi_k\pi_j})
         + (c_{\pi_k\pi_i} - c_{\pi_i\pi_k})\bigr].
    \label{eq:delta_exchange}
\end{equation}
Each evaluation costs $O(|i - j|)$, and a full scan costs $O(n^3)$ without
further optimisation.

\paragraph{Insert.} Removes element $\pi_i$ from its position and reinserts
it at position $j \neq i$. The neighbourhood has $(n-1)^2$ neighbours. As
shown in~\cite{SchStu04}, any insert move is equivalent to a sequence of
$|i - j|$ adjacent transpose moves, which allows incremental evaluation:
the delta for inserting $\pi_i$ at position $j$ is accumulated from the
intermediate swap deltas. This brings the cost of a full scan down from
$O(n^3)$ to $O(n^2)$.

\section{Initialisation}

\paragraph{Random permutation.} A uniformly random permutation is generated
using a Fisher-Yates shuffle. It is simple and covers different regions of the
search space each time.

\paragraph{Chenery-Watanabe (CW) heuristic.} Originally introduced for
triangularising economic input-output tables~\cite{CheWat58}, this greedy
method builds a permutation by repeatedly selecting the element $i$ not yet
placed that has the largest row sum:
\begin{equation}
    a_i = \sum_{j=1}^{n} c_{ij}.
    \label{eq:cw}
\end{equation}
The element with the highest $a_i$ is ranked first, and the process continues
until all elements are placed. CW runs in $O(n^2)$ and typically produces a
much better starting point than a random permutation.

\section{Iterative Improvement}

An iterative improvement algorithm starts from an initial solution and
repeatedly applies an improving move until no such move exists. It stops at a
local optimum. Two pivoting strategies were used.

\paragraph{First-improvement (FI).} The neighbourhood is scanned in a fixed
order and the first move with a strictly positive delta is applied immediately.
This avoids scanning the full neighbourhood at each step.

\paragraph{Best-improvement (BI).} The entire neighbourhood is scanned and the
move with the highest delta is applied. Each iteration is more expensive but
makes a larger improvement per step.

All six combinations of pivoting rule and neighbourhood were tested with both
initialisation strategies, giving $2 \times 3 \times 2 = 12$ configurations.

\section{Variable Neighbourhood Descent}

Variable Neighbourhood Descent (VND) runs a sequence of local searches, each
using a different neighbourhood~\cite{HanMla01}. If a search in neighbourhood
$\mathcal{N}_k$ finds an improvement, the algorithm restarts from
$\mathcal{N}_1$. If no improvement is found, it moves to $\mathcal{N}_{k+1}$.
The process stops when no neighbourhood can improve the solution.

Two orderings were tested, both using first-improvement and CW initialisation:
\begin{itemize}
    \item \textbf{VND1:} Transpose $\to$ Exchange $\to$ Insert.
    \item \textbf{VND2:} Transpose $\to$ Insert $\to$ Exchange.
\end{itemize}
Transpose is placed first in both cases because it is the cheapest
neighbourhood to scan.

\section{Experimental Setup}

All algorithms were implemented in C and compiled with \texttt{-O3}. Each
configuration was applied once to each of the 78 benchmark instances. Solution
quality is measured using the relative percentage deviation (RPD) from the
best-known value $f^*$:
\begin{equation}
    \text{RPD} = 100 \times \frac{f^* - f(\pi)}{f^*}.
    \label{eq:rpd}
\end{equation}
An RPD of 0 means the best-known solution was reached; lower is better.
Note that this convention gives non-negative values, since $f(\pi) \leq f^*$
by definition of the best-known solution. Some references use the opposite
sign convention, which produces non-positive values, but the interpretation
is equivalent.

Statistical comparisons between algorithm pairs use a paired Wilcoxon
signed-rank test and a paired t-test at significance level $\alpha = 0.05$,
run in R. Both tests are paired because each algorithm is evaluated on the
same 78 instances, which removes the effect of instance difficulty. The
Wilcoxon test makes no normality assumption; the t-test is included as a
standard complement.
